\begin{abstract}
Gate-count optimality for sparse quantum state preparation does not by itself imply depth optimality, particularly when no qubits beyond the final output register are available. We consider the preparation of an $n$-qubit state with $d$ nonzero complex amplitudes from an explicit classical description, using arbitrary single-qubit unitaries and CNOT gates in an all-to-all, purely unitary model. We present a constructive argument that, for $n\ge2$ and $2\le d\le2^n$, every target admits a circuit on exactly the $n$ output wires with simultaneous depth $O(\log(n+1)+d/\log(n+1))$ and gate count $O(n+nd/\log(n+1))$. For each fixed $c\ge1$, these bounds match the corresponding worst-case lower bounds in the window $d\le\min\{n^c,2^n\}$; no optimality claim is made for general exponentially large support. The construction combines width-constrained nonlinear feature evaluation, recoverable workspace within the output register, and a batched reduction that turns small-support preparation circuits into full-space Householder reflections.

At the implementation level, we give an exact reflection-selection rule that preserves both resource bounds, and separately certify rotation-rounding errors, repeated-reflection errors, and classical preprocessing residuals. For explicitly specified complex rational inputs, an a posteriori interval verifier encloses the source circuit's output without constructing a full $2^n$-dimensional state vector. Combining this enclosure with a spectral-norm compilation bound yields a target-state error guarantee for accepted instances. Finite tests report resource changes within the same prototype and the rejection of invalid certificates; they do not establish a hardware advantage or Clifford+$T$ optimality.
\end{abstract}
\noindent\textbf{Keywords:} sparse state preparation; quantum circuit synthesis; depth complexity; ancillas; Householder reflections; error certification.

\smallskip
{\small\draftstatus}
